\chapter{เวกเตอร์ในทางฟิสิกส์และการวิเคราะห์แรง}
\label{ch:ch02_vectors}
\index{เวกเตอร์และการรวมแรง}

\begin{conceptbox}[title=วัตถุประสงค์การเรียนรู้ประจำบท]
เมื่อสิ้นสุดการศึกษาบทเรียนนี้ นักศึกษาสามารถ
\begin{enumerate}
    \item เมื่อศึกษาบทนี้จบแล้ว ผู้เรียนสามารถ
    \item แยกความแตกต่างระหว่างปริมาณสเกลาร์และปริมาณเวกเตอร์ พร้อมยกตัวอย่างปริมาณทางฟิสิกส์แต่ละประเภทได้ถูกต้อง
    \item บวกและลบเวกเตอร์ด้วยวิธีเรขาคณิต ทั้งวิธีหัวต่อหางและวิธีสี่เหลี่ยมด้านขนาน
    \item แตกเวกเตอร์ออกเป็นองค์ประกอบตามแกน x และ y และคำนวณกลับหาขนาดและทิศทางของเวกเตอร์ลัพธ์ด้วยวิธีองค์ประกอบ
    \item คำนวณผลคูณเชิงสเกลาร์ (dot product) และอธิบายความหมายทางฟิสิกส์ เช่น การหางาน
\end{enumerate}
\end{conceptbox}

\section{ทฤษฎีและหลักการพื้นฐาน}

\begin{bioappbox}
\textbf{การวิเคราะห์เวกเตอร์ของแรงและความเร็วในกระบวนการทางชีวภาพ}

การเคลื่อนที่ของสารละลายผ่านเยื่อหุ้มเซลล์และการเคลื่อนที่ของสิ่งมีชีวิตเซลล์เดียว เช่น \textit{Paramecium} ที่ว่ายน้ำด้วยการโบกพัดของซิเลีย (Cilia) สามารถจำลองได้ด้วยเวกเตอร์ความเร็วสองมิติและสามมิติ นอกจากนี้ แรงตึงผิวและแรงดึงดูดของของเหลวในท่อลำเลียงไซเลม (Xylem) ของพืชยังต้องวิเคราะห์ด้วยการแตกแรงเวกเตอร์ในแนวแกนดิ่งและแกนราบ เพื่ออธิบายกลไกการดูดน้ำจากรากขึ้นสู่ยอดไม้สูงกว่า 20 เมตรได้อย่างสอดคล้องกับหลักกลศาสตร์
\end{bioappbox}

\begin{labbox}
1. ตรวจสอบตำแหน่งศูนย์ (Zero Error) ของเครื่องมือวัดทุกชนิดก่อนเริ่มบันทึกข้อมูล
2. ทำการวัดซ้ำอย่างน้อย 3 ถึง 5 ครั้งในแต่ละจุดทดลองเพื่อคำนวณหาค่าเฉลี่ยและค่าเบี่ยงเบนมาตรฐาน
3. บันทึกผลการวัดพร้อมระบุหน่วยเอสไอ (SI Units) และค่าความคลาดเคลื่อนของการวัดเสมอ
\end{labbox}

\section{ตัวอย่างโจทย์และการวิเคราะห์เชิงฟิสิกส์}

\begin{examplebox}
ตัวอย่างที่ 2.1 เวกเตอร์แรง A มีขนาด 20 นิวตัน ทำมุม 30$^\circ$ กับแกน x บวก และเวกเตอร์แรง B มีขนาด 15 นิวตัน ทำมุม 120$^\circ$ กับแกน x บวก จงหาเวกเตอร์แรงลัพธ์ด้วยวิธีองค์ประกอบ

วิธีทำ แตกองค์ประกอบของแต่ละเวกเตอร์ A$_x$ = 20 cos30$^\circ$ = 20 $\times$ 0.8660 = 17.32 N A$_y$ = 20 sin30$^\circ$ = 20 $\times$ 0.5000 = 10.00 N B$_x$ = 15 cos120$^\circ$ = 15 $\times$ (−0.5000) = −7.50 N B$_y$ = 15 sin120$^\circ$ = 15 $\times$ 0.8660 = 12.99 N

รวมองค์ประกอบ R$_x$ = 17.32 + (−7.50) = 9.82 N R$_y$ = 10.00 + 12.99 = 22.99 N

หาขนาดและทิศทาง R = √(9.82$^2$ + 22.99$^2$) = √(96.4 + 528.5) = √624.9 $\approx$ 25.0 N $\theta$ = tan$^-$$^1$(22.99/9.82) = tan$^-$$^1$(2.341) $\approx$ 66.9$^\circ$
\end{examplebox}

\section{แบบฝึกหัดทบทวนและคำถามท้ายบท}

\begin{enumerate}
    \item จงบอกว่าปริมาณต่อไปนี้เป็นปริมาณสเกลาร์หรือปริมาณเวกเตอร์: อุณหภูมิ ความเร่ง พลังงาน โมเมนตัม เวลา และแรง
    \item เวกเตอร์ A มีขนาด 12 หน่วย ทำมุม 40$^\circ$ กับแกน x บวก จงหาองค์ประกอบ A$_x$ และ A$_y$
    \item เวกเตอร์ B มีองค์ประกอบ B$_x$ = −5 และ B$_y$ = 9 จงหาขนาดและทิศทางของเวกเตอร์ B
    \item จงหาผลบวกของเวกเตอร์ C (ขนาด 8 หน่วย ทำมุม 0$^\circ$) และเวกเตอร์ D (ขนาด 6 หน่วย ทำมุม 90$^\circ$) ด้วยวิธีองค์ประกอบ พร้อมหาขนาดและทิศทางของเวกเตอร์ลัพธ์
    \item เรือลำหนึ่งแล่นด้วยความเร็ว 10 กิโลเมตรต่อชั่วโมงไปทางทิศเหนือ ขณะที่กระแสน้ำไหลด้วยความเร็ว 4 กิโลเมตรต่อชั่วโมงไปทางทิศตะวันออก จงหาขนาดและทิศทางของความเร็วลัพธ์ของเรือ
    \item เวกเตอร์แรง F₁ = (3, 4) นิวตัน และ F₂ = (−1, 2) นิวตัน จงหาแรงลัพธ์ F₁ + F₂ และขนาดของแรงลัพธ์นั้น
\end{enumerate}

% สิ้นสุดบทเรียน
